\documentclass[11pt]{article}
\usepackage[margin=1in]{geometry}
\usepackage{amsmath, amssymb}
\usepackage{xcolor}
\usepackage{mdframed}
\usepackage{parskip}
\usepackage{enumitem}

\newmdenv[
  linecolor=black!25,
  backgroundcolor=black!3,
  roundcorner=4pt,
  skipabove=8pt,
  skipbelow=8pt,
  innertopmargin=6pt,
  innerbottommargin=6pt
]{keybox}

\title{Sub-pixel Extrema Refinement and Edge Rejection in SIFT}
\author{}
\date{}

\begin{document}
\maketitle

\section{The one-dimensional skeleton}

Expand $f$ around a sample point $a$, and write the offset as $\delta = x - a$:
\[
f(a + \delta) \;\approx\; f(a) + f'(a)\,\delta + \tfrac{1}{2} f''(a)\,\delta^2 .
\]
To locate the critical point, differentiate \emph{with respect to $\delta$} and set the
result to zero:
\[
\frac{d}{d\delta}\Big[\, f(a) + f'(a)\,\delta + \tfrac{1}{2}f''(a)\,\delta^2 \,\Big]
= f'(a) + f''(a)\,\delta = 0 .
\]
\begin{keybox}
\[
\delta = -\frac{f'(a)}{f''(a)} .
\]
\end{keybox}
The refined location is $a + \delta$. This is exactly one Newton step applied to $f'$,
which is what setting a quadratic's derivative to zero amounts to.

\section{Lifting to three dimensions}

Now the coordinate is $\mathbf{x} = (s, y, x)$ and $f = D$ is the DoG value. The single
derivative becomes a \textbf{gradient} (a $3$-vector); the second derivative becomes a
\textbf{Hessian} (a $3\times 3$ matrix):
\[
\mathbf{g} = \nabla D =
\begin{bmatrix} D_s \\ D_y \\ D_x \end{bmatrix},
\qquad
H =
\begin{bmatrix}
D_{ss} & D_{sy} & D_{sx} \\
D_{sy} & D_{yy} & D_{yx} \\
D_{sx} & D_{yx} & D_{xx}
\end{bmatrix}.
\]
With offset vector $\boldsymbol{\delta} = (\delta_s, \delta_y, \delta_x)$, the expansion is
\[
D(\mathbf{x}_0 + \boldsymbol{\delta}) \;\approx\;
D(\mathbf{x}_0) + \mathbf{g}^{\top}\boldsymbol{\delta}
+ \tfrac{1}{2}\,\boldsymbol{\delta}^{\top} H \boldsymbol{\delta}.
\]
This is the $1$D formula with each piece promoted:
$f'(a)\,\delta \to \mathbf{g}^{\top}\boldsymbol{\delta}$ (a dot product), and
$\tfrac{1}{2} f''(a)\,\delta^2 \to \tfrac{1}{2}\boldsymbol{\delta}^{\top} H\boldsymbol{\delta}$
(a quadratic form).

\section{Solving for the offset}

Take the gradient with respect to $\boldsymbol{\delta}$ and set it to $\mathbf{0}$. Two
standard identities do the work --- the multivariable analogues of
$\frac{d}{d\delta}(c\,\delta) = c$ and
$\frac{d}{d\delta}(\tfrac{1}{2} q\,\delta^2) = q\,\delta$:
\[
\frac{\partial}{\partial \boldsymbol{\delta}}\big(\mathbf{g}^{\top}\boldsymbol{\delta}\big)
= \mathbf{g},
\qquad
\frac{\partial}{\partial \boldsymbol{\delta}}
\big(\tfrac{1}{2}\boldsymbol{\delta}^{\top} H \boldsymbol{\delta}\big)
= H\boldsymbol{\delta}
\quad (\text{$H$ symmetric}).
\]
Therefore
\[
\frac{\partial D}{\partial \boldsymbol{\delta}}
= \mathbf{g} + H\boldsymbol{\delta} = \mathbf{0},
\]
\begin{keybox}
\[
H\boldsymbol{\delta} = -\mathbf{g}
\qquad\Longrightarrow\qquad
\boldsymbol{\delta} = -H^{-1}\mathbf{g}.
\]
\end{keybox}
Compare the $1$D box, $\delta = -f'/f''$: identical structure. Here $-H^{-1}$ \emph{is}
``divide by the second derivative,'' except that in three dimensions ``divide'' means
``solve a $3\times 3$ linear system.''

\section{What you do with the offset}

\subsection*{1. Refined location}
\[
\hat{\mathbf{x}} = \mathbf{x}_0 + \boldsymbol{\delta}.
\]
Because $s$ is one of the three axes, the scale is refined too --- $\sigma$ becomes
continuous instead of snapping to the discrete level ladder.

\subsection*{2. Refined value at the peak}
Substitute $\boldsymbol{\delta} = -H^{-1}\mathbf{g}$ back into the expansion. The quadratic
term simplifies using $H\boldsymbol{\delta} = -\mathbf{g}$:
\[
\boldsymbol{\delta}^{\top} H \boldsymbol{\delta}
= \boldsymbol{\delta}^{\top}(H\boldsymbol{\delta})
= \boldsymbol{\delta}^{\top}(-\mathbf{g})
= -\,\mathbf{g}^{\top}\boldsymbol{\delta}.
\]
Hence
\[
D(\hat{\mathbf{x}})
= D + \mathbf{g}^{\top}\boldsymbol{\delta}
+ \tfrac{1}{2}\boldsymbol{\delta}^{\top} H \boldsymbol{\delta}
= D + \mathbf{g}^{\top}\boldsymbol{\delta} - \tfrac{1}{2}\mathbf{g}^{\top}\boldsymbol{\delta},
\]
\begin{keybox}
\[
D(\hat{\mathbf{x}}) = D(\mathbf{x}_0) + \tfrac{1}{2}\,\mathbf{g}^{\top}\boldsymbol{\delta}.
\]
\end{keybox}
The peak of a parabola sits half a gradient-step above the sample --- the same $\tfrac12$
as in $1$D. Threshold $|D(\hat{\mathbf{x}})| < 0.03$ to reject low-contrast points.

\subsection*{3. Convergence / re-centering}
If $\max_i |\delta_i| > 0.5$, the true peak is nearer a different sample. Move
$\mathbf{x}_0$ to that integer neighbour, recompute $\mathbf{g}$ and $H$ there, and re-solve.
Iterate at most $5$ times; discard points that fail to converge or wander outside the volume.
Every derivative is thus always evaluated at an integer grid point --- only the final
residual offset is fractional.

\section{Finite-difference stencils}

The gradient and Hessian are filled in by central differences on the integer grid
(spacing $1$):
\[
D_x = \frac{D[x{+}1] - D[x{-}1]}{2},
\qquad
D_{xx} = D[x{+}1] - 2D[x] + D[x{-}1],
\]
\[
D_{xy} = \frac{D[y{+}1,x{+}1] - D[y{+}1,x{-}1] - D[y{-}1,x{+}1] + D[y{-}1,x{-}1]}{4}
\]
(indices along the other axes suppressed). The mixed-partial stencil touches all corners
of the $3\times 3\times 3$ cube --- which is why the border voxels must be dropped before
this step.

\section{Edge rejection via the Hessian}

DoG responds strongly all along an edge, producing extrema strung out down a ridge that
slide under small viewpoint changes. The discriminator is curvature. Take the
\emph{spatial} $2\times 2$ Hessian (drop the scale row and column):
\[
H_{2\times 2} =
\begin{bmatrix} D_{xx} & D_{xy} \\ D_{xy} & D_{yy} \end{bmatrix}.
\]
Its eigenvalues $\alpha, \beta$ are the principal curvatures. A ridge has one large and one
small; a blob has both large. Their ratio $r = \alpha/\beta$ can be tested without an
eigendecomposition, using only trace and determinant:
\[
\frac{\operatorname{Tr}(H_{2\times 2})^2}{\det(H_{2\times 2})}
= \frac{(\alpha + \beta)^2}{\alpha\beta}
= \frac{(r + 1)^2}{r}.
\]
This depends only on the ratio, not on the individual magnitudes, and is minimised at
$r = 1$. Keep a point when
\begin{keybox}
\[
\frac{\operatorname{Tr}(H_{2\times 2})^2}{\det(H_{2\times 2})} < \frac{(r+1)^2}{r},
\qquad r = 10
\quad\Big(\text{threshold } = \tfrac{121}{10} = 12.1\Big).
\]
\end{keybox}
Reject $\det(H_{2\times 2}) \le 0$ outright: opposite-sign curvatures mean a saddle, not an
extremum, and guard against dividing by near-zero.

\section*{Pipeline order}

Refine first, then test contrast on the refined value $D(\hat{\mathbf{x}})$, then apply edge
rejection. Doing edge rejection before refinement tests the Hessian at the wrong location.

\end{document}
